\begin{pastexam}[label={px:ccu-110-swa-07}]{中正資工所甲組 110 學年度 第 9 題}
（證明）

(3\%) Prove or disprove with a counter example that every comparison based sorting algorithm has
time complexity $\Omega(n\lg n)$.
\end{pastexam}

\begin{pastsoln}
題目沒有說是哪一種情況的時間。 「時間複雜度」沒有特別說明時指最壞情況，這裡照這個解讀：
敘述成立，下面證明。 最後再說明換成最好情況時為什麼不成立。

\textbf{證明。} 設 $\mathcal{A}$ 是任何一個正確的比較排序，固定 $n$。 只看 $n$ 個鍵兩兩不同的
輸入：這些輸入已經需要那麼多次比較，全部輸入的最壞情況只會更多。 鍵兩兩不同時，每次比較只有
$\le$、$>$ 兩種結果。
\begin{enumerate}
\item 把 $\mathcal{A}$ 對這些輸入的所有執行過程畫成決策樹：內部節點是一次比較 $i{:}j$，左右子樹
對應 $a_i \le a_j$ 和 $a_i > a_j$；葉節點是演算法停下來的地方。 一次執行的比較次數就是走到的那個
葉節點的深度。
\item 葉節點至少 $n!$ 個。 $n$ 個相異的鍵有 $n!$ 種大小順序。 假設兩種不同的順序 $\pi \ne \sigma$
走到同一個葉節點，$\mathcal{A}$ 在兩個輸入上看到的比較結果完全一樣，就會做完全一樣的搬動，輸出
同一種重排；同一種重排不可能同時把兩種不同的順序都排好，其中一個會排錯，和 $\mathcal{A}$ 正確
矛盾。 所以 $n!$ 種順序走到 $n!$ 個不同的葉節點。
\item 深度都不超過 $d$ 的二元樹最多有 $2^d$ 個葉節點。 對 $d$ 歸納：$d = 0$ 時只有根，$1$ 個葉
節點；$d \ge 1$ 時左右子樹的葉節點深度都不超過 $d - 1$，各自最多 $2^{d-1}$ 個，合起來最多
$2^d$ 個。
\item 設最壞情況比較 $d$ 次，也就是最深的葉節點深度是 $d$。 由 2、3，$2^d \ge n!$，所以
$d \ge \lg(n!)$。 $n!$ 裡最大的 $\lceil n/2 \rceil$ 個因數都不小於 $n/2$，所以
$n! \ge (n/2)^{n/2}$，$\lg(n!) \ge \frac{n}{2}(\lg n - 1)$。 $n \ge 4$ 時
$\lg n - 1 \ge \frac{1}{2}\lg n$，得到 $d \ge \frac{n}{4}\lg n$。
\item 每次比較至少花常數時間，所以 $\mathcal{A}$ 的最壞執行時間至少是常數乘上 $\frac{n}{4}\lg n$，
也就是 $\Omega(n \lg n)$。
\end{enumerate}
$\mathcal{A}$ 是任意的，所以每一個比較排序的最壞時間都是 $\Omega(n \lg n)$。 這就是資結書
定理~15.1；資結書定理~15.3 說平均情況也是 $\Omega(n \lg n)$。

\textbf{換成最好情況時不成立。} 反例是插入排序（資結書第~1.3~節）：輸入已經由小到大排好時，外層
迴圈每一輪只比一次、不挪動，共 $n - 1$ 次比較，時間 $\Theta(n)$。 所以「每個比較排序的
\textbf{最好}情況都是 $\Omega(n \lg n)$」是錯的。 作答時兩種情況都寫出來最安全。
\end{pastsoln}
